% Options for packages loaded elsewhere
\PassOptionsToPackage{unicode}{hyperref}
\PassOptionsToPackage{hyphens}{url}
\documentclass[10pt]{article}
\usepackage[letterpaper,margin=1in]{geometry}
\usepackage{fontspec}
\setmainfont{Arial}
\usepackage{graphicx}
\usepackage{float}
\usepackage{setspace}
\usepackage{lineno}
\doublespacing
\linenumbers
\usepackage{xcolor}
\usepackage{amsmath,amssymb}
\setcounter{secnumdepth}{-\maxdimen} % remove section numbering
\usepackage{iftex}
\ifPDFTeX
  \usepackage[T1]{fontenc}
  \usepackage[utf8]{inputenc}
  \usepackage{textcomp} % provide euro and other symbols
\else % if luatex or xetex
  \usepackage{unicode-math} % this also loads fontspec
  \defaultfontfeatures{Scale=MatchLowercase}
  \defaultfontfeatures[\rmfamily]{Ligatures=TeX,Scale=1}
\fi
\usepackage{lmodern}
\ifPDFTeX\else
  % xetex/luatex font selection
\fi
% Use upquote if available, for straight quotes in verbatim environments
\IfFileExists{upquote.sty}{\usepackage{upquote}}{}
\IfFileExists{microtype.sty}{% use microtype if available
  \usepackage[]{microtype}
  \UseMicrotypeSet[protrusion]{basicmath} % disable protrusion for tt fonts
}{}
\makeatletter
\@ifundefined{KOMAClassName}{% if non-KOMA class
  \IfFileExists{parskip.sty}{%
    \usepackage{parskip}
  }{% else
    \setlength{\parindent}{0pt}
    \setlength{\parskip}{6pt plus 2pt minus 1pt}}
}{% if KOMA class
  \KOMAoptions{parskip=half}}
\makeatother
\usepackage{longtable,booktabs,array}
\newcounter{none} % for unnumbered tables
\usepackage{calc} % for calculating minipage widths
% Correct order of tables after \paragraph or \subparagraph
\usepackage{etoolbox}
\makeatletter
\patchcmd\longtable{\par}{\if@noskipsec\mbox{}\fi\par}{}{}
\makeatother
% Allow footnotes in longtable head/foot
\IfFileExists{footnotehyper.sty}{\usepackage{footnotehyper}}{\usepackage{footnote}}
\makesavenoteenv{longtable}
\setlength{\emergencystretch}{3em} % prevent overfull lines
\providecommand{\tightlist}{%
  \setlength{\itemsep}{0pt}\setlength{\parskip}{0pt}}
\usepackage{bookmark}
\IfFileExists{xurl.sty}{\usepackage{xurl}}{} % add URL line breaks if available
\urlstyle{same}
% fallback for those not using the hyperref driver hyperxmp:
\makeatletter
\@ifundefined{xmpquote}{\newcommand{\xmpquote}[1]{#1}}{}
\makeatother
\hypersetup{
  hidelinks,
  pdfcreator={LaTeX via pandoc}}

\author{}
\date{}

\begin{document}

\begin{center}
{\LARGE\bfseries Benchmark Hygiene and the Non-Canonical Frontier in Peptide-MHC-I Binding Prediction\par}
\vspace{1em}
{\large Yunshi Huang$^{1,\dagger}$, Weisheng Chen$^{2,\dagger}$, Yiqi Cai$^{2}$, Xiaopeng Zhuang$^{2}$, Zhuxian Zhou$^{3}$, Xiangyang Xue$^{4,*}$, Zhuha Zhou$^{2,*}$\par}
\vspace{0.6em}
{\normalsize $^{1}$Department of General Surgery, The First Affiliated Hospital of Wenzhou Medical University, Wenzhou, China\par}
{\normalsize $^{2}$Department of Gastroenterology Surgery, The First Affiliated Hospital of Wenzhou Medical University, Wenzhou, China\par}
{\normalsize $^{3}$Center for Bionanoengineering, College of Chemical and Biological Engineering, Zhejiang University, Hangzhou, China\par}
{\normalsize $^{4}$Zhejiang Key Laboratory of Intelligent Cancer Biomarker Discovery and Translation, The First Affiliated Hospital of Wenzhou Medical University, Wenzhou, China\par}
\vspace{0.6em}
{\normalsize $\dagger$ These authors contributed equally to this work.\par}
{\normalsize $*$ Corresponding authors: wzxxy@wmu.edu.cn (XX); zhouzhuha@wmu.edu.cn (ZZ)\par}
\end{center}

\subsection{Abstract}\label{abstract}

\textbf{Motivation}: Peptide-MHC-I binding-prediction benchmarks can
manufacture false conclusions when confounders are uncontrolled.
\textbf{Results}: We benchmarked six architectures on 5,190
MHCflurry-labelled HLA-A*02:01 9-mers with multi-seed estimates and
re-analysed three apparent findings under controlled conditions. A
convolutional network was modestly best (91.9\% $\pm$ 0.9\% vs 90.6\% $\pm$
1.2\% for a deep feed-forward network); random labels fell to chance
(\textasciitilde33\%); and a claimed `label quality dominates
architecture' result dissolved once the degenerate null control was
removed. A zero-parameter anchor rule matched the trained model on
canonical peptides, while non-canonical binders plateaued at
\textasciitilde0.65 sensitivity across three independent predictors ---
the true frontier. Extending to five common HLA-I alleles and a
15-allele panel (14,488 curated IEDB positives), a rebalanced
pan-specific model matched per-allele specialists on four of five
alleles and reached 0.819 pooled non-canonical sensitivity. A temporal
holdout of post-2022 epitopes preserved recognition, and a benchmark of
6,524 experimentally measured affinities (ROC AUC 0.942) confirmed the
training labels' quality. We distil the artifacts into a
benchmark-hygiene checklist applicable beyond peptide-MHC.
\textbf{Availability}: Code, models, and data at
github.com/woodhaha/peptide-mhc-benchmark. \textbf{Contact}:
zhouzhuha@wmu.edu.cn.

\begin{center}\rule{0.5\linewidth}{0.5pt}\end{center}

\subsection{Author Summary}\label{author-summary}

Our immune system detects infected and cancerous cells by reading short protein fragments, called peptides, displayed on the cell surface. Computer programs that predict which peptides will be displayed are central to designing vaccines and cancer immunotherapies. Here we show that the standard way of testing these programs can make them appear more accurate than they really are --- a subtle circularity in which a program is evaluated against data it helped generate. By re-running our own benchmarks under careful controls, we uncovered three ways this circularity produces misleading results. We then asked which peptides remain genuinely hard to predict. Peptides that follow the usual rules are now predicted almost perfectly, but a minority that breaks those rules remains a persistent blind spot for every method we tested. We distill our findings into a practical checklist to help researchers build more trustworthy prediction benchmarks, and we release all code, models, and data so the work can be verified and extended.

\subsection{1 Introduction}\label{introduction}

MHC class I molecules present 8--11-mer peptides to CD8+ T cells, and
accurate prediction of peptide--MHC-I binding underpins epitope
discovery for vaccine design and cancer immunotherapy (Falk et al.,
1991; Lundegaard et al., 2011). For HLA-A*02:01 --- the most prevalent
class I allele globally --- binding is dominated by two primary anchor
positions: position 2 (p2, hydrophobic: Leu/Met/Ile/Val) and position 9
(p9, C-terminal, Val/Leu preferred) (Ruppert et al., 1993; Rammensee et
al., 1995). Computational predictors have evolved from position-specific
scoring matrices to neural networks and pan-specific tools such as
NetMHCpan \cite{Reynisson20}, MHCflurry \cite{ODonnell20},
BigMHC \cite{Albert23}, and MUNIS \cite{Wohlwend25},
increasingly using protein language model embeddings (Rives et al.,
2021).

Several gaps remain. Head-to-head architecture comparisons under fully
controlled conditions (identical training data, encoding, evaluation)
are uncommon. The impact of training-label quality on downstream
performance is seldom quantified, despite evidence that \textgreater55\%
of IEDB entries are computationally rather than experimentally labelled
\cite{Preibisch26}. Whether a single pan-specific model can replace per-allele specialists, and how hard non-canonical binders are across alleles, remains unsettled. And the translation of binding prediction to
real-world epitope discovery is rarely demonstrated in a unified,
reproducible pipeline.

We address these gaps through a systematic benchmark designed to isolate
the effects of architecture choice, label quality, and encoding. Our
contributions are: (i) a controlled six-architecture comparison with
multi-seed uncertainty estimates, showing a modest but measurable
architecture effect; (ii) a label-source control confirming that random
labels fall to chance; (iii) evaluation of ESM-2 embeddings against
BLOSUM62 with a mean-pooled ablation; (iv) held-out validation against
49 IEDB epitopes (44 of which overlap the labeler's training corpus);
(v) translational application through protein scanning and cancer
mutation analysis; and (vi) a multi-allele extension across five common
HLA-I alleles and a 15-allele panel, testing whether a single rebalanced
pan-specific model matches per-allele specialists and quantifying how
hard non-canonical binders are for three independent sequence
predictors.

\subsection{2 Methods}\label{methods}

\subsubsection{2.1 Training data and
labels}\label{training-data-and-labels}

100,000 random 9-mer peptides were generated by uniform sampling of the
20 standard amino acids. Each was labelled with MHCflurry 2.2.1 against
HLA-A*02:01 and assigned to Strong Binder (SB, percentile rank
\textless0.5\%), Weak Binder (WB, 0.5\% $\leq$ rank\textless2.0\%), or
Non-Binder (NB, rank  $\geq$ 2.0\%) following the netMHCpan convention
\cite{Reynisson20}. The resulting distribution (1.7\% SB, 5.1\%
WB, 93.2\% NB) was balanced by downsampling to 1,730 per class
(n=5,190), partitioned 90/10 into training (n=4,671) and held-out test
(n=519) sets with stratified sampling (seed 42). Two comparator
labelling strategies were implemented: a PSSM encoding crystallographic
A*02:01 anchor preferences, and random synthetic labels as a baseline
control. A held-out benchmark of 49 experimentally validated IEDB
epitopes (HLA-A*02:01-restricted, with positive binding/T-cell assay
evidence) plus 20 homopolymer negatives ( $\leq$ 2 unique amino acids) was
curated from the IEDB (query-api.iedb.org, accessed 2026-08-02);
label-circularity overlap (44/49) was determined by sequence-identity
match against the IEDB entries in MHCflurry's training corpus.

\subsubsection{2.2 Peptide encoding}\label{peptide-encoding}

Each 9-mer was encoded with the BLOSUM62 matrix (Henikoff and Henikoff,
1992), min--max normalized, producing a 9$\times$20 matrix flattened to 180
dimensions for fully connected networks. ESM-2 embeddings (Lin et al.,
2023) were extracted at t6 (2,880-dim per-peptide; 320 per position $\times$ 9)
and t12 (4,320-dim; 480 per position $\times$ 9) depths --- the two smallest
ESM-2 checkpoints (8M and 35M parameters); a mean-pooled t6 variant
(320-dim) served as an ablation control destroying positional
information.

\subsubsection{2.3 Model architectures}\label{model-architectures}

Six classifiers were implemented: (i) FFN \cite{Nielsen03}:
180 $
ightarrow$ 90 $
ightarrow$ 3, 49K params; (ii) Deep FFN: 180 $
ightarrow$ 360 $
ightarrow$ 180 $
ightarrow$ 90 $
ightarrow$ 45 $
ightarrow$ 3 with batch
normalisation and dropout, 152K params; (iii) CNN: Conv2D(32,3$\times$3) + FFN
body, 1.05M params; (iv) LSTM(64) $
ightarrow$ dense(32) $
ightarrow$ 3, 24K params; (v)
ResNet-style stem + 3 residual blocks, 325K params; (vi) Random Forest
(100 trees). For ESM-2 embeddings an FFN classifier with a 256 $
ightarrow$ 128 $
ightarrow$ 64
stack was trained. All networks used cross-entropy loss with Adam or
RMSprop (lr 0.001), 150 epochs, batch size 50, 20\% internal validation
split, early stopping (patience 10), and learning-rate reduction on
plateau. Five-fold stratified cross-validation was performed for the
FFN. The held-out test set (519 peptides) was disjoint from training and
used for all reported metrics.

\subsubsection{2.4 Protein scanning, mutation analysis, and
structure}\label{protein-scanning-mutation-analysis-and-structure}

Ten therapeutically relevant proteins (MART-1, gp100, tyrosinase,
NY-ESO-1, WT1, p53, KRAS, influenza M1, CMV pp65, SARS-CoV-2 spike RBD)
were scanned with a sliding 9-mer window (stride 1; 3,536 windows).
Sixteen recurrent hotspot mutations (7 p53, 9 KRAS) were analyzed by
comparing wild-type and mutant binding scores; mutations were classified
as Created/Enhanced/Destroyed/Unchanged. The KRAS G12V neoepitope
YKLVVVGAV carries a non-canonical lysine at P2.

\subsubsection{2.5 Feature extraction}\label{feature-extraction}

Learned features were extracted from the penultimate dense layer via
truncated models. Hierarchical clustering (Ward's method, Euclidean
distance) was applied to representations of top epitope candidates;
feature--binding-score correlation was assessed by Pearson's r with a
1,000-iteration permutation test to control for chance correlations.
t-SNE visualisation of the 45-dimensional penultimate layer was used to
inspect class separation, and integrated gradients identified
position-specific anchor contributions.

\subsubsection{2.6 Multi-allele extension}\label{multi-allele-extension}

\paragraph{2.6.1 Multi-allele data
curation}\label{multi-allele-data-curation}

Experimentally confirmed class I binders were curated from the IEDB
\cite{Vita19} for five common alleles (HLA-A*02:01, A*01:01,
A*03:01, B*07:02, B*44:03), yielding n = 14,488 positives (canonical
11,225; Type A 151; Type B 3,112). A 15-allele panel (pan16: 11
additional alleles --- A*11:01, A*24:02, B*08:01, B*15:01, B*27:05,
B*35:01, B*40:01, B*49:01, B*51:01, B*57:01) was curated for
generalisation testing. Protein-background negatives were sampled from
natural proteomes to avoid the uniform-random decoys used in earlier
evaluations. IEDB data were queried live from query-api.iedb.org
(accessed 2026-08-02); the 15-allele and temporal-holdout estimates are
therefore snapshot-dependent and drift slightly as the IEDB grows.

\paragraph{2.6.2 Type A/B classification}\label{type-ab-classification}

Type A (tolerated weak anchors) and Type B (energy-compensated
non-canonical) were defined by whether the peptide violates the allele's
canonical anchor motif at p2/p9. Type A carries broad-only tolerated
anchors (residues that appear in binding motifs but are non-canonical);
Type B carries non-canonical residues at anchor positions. The two
classes have different biological origins and, as we show, different
prediction difficulty.

\paragraph{2.6.3 Pan model and per-allele
comparators}\label{pan-model-and-per-allele-comparators}

The pan-specific DFFN uses BLOSUM62 encoding (Henikoff and Henikoff,
1992) through a 360--180--90--45 hidden stack with batch normalisation
and dropout. Rebalancing upsampled the minority B*07:02 allele 4$\times$ during
training to correct for allele-frequency imbalance. Five seeds (0, 1,
42, 123, 2024) were used; per-seed variation arises from data-sampling
RNG. Per-allele models (one per allele) and a pan16 model (15 alleles,
per-allele minimum 2,000 upsampling) were trained with the same recipe.

\paragraph{2.6.4 Comparator predictors}\label{comparator-predictors}

MHCflurry 2.2.1 \cite{ODonnell20} served as an independent
sequence predictor (affinity threshold 500 nM). MUNIS (Wohlwend et al.,
2025) --- a second sequence-based predictor using ESM-2 embeddings
\cite{Rives21} --- served as an independent sequence comparator,
run on the n = 3,112 Type B population plus canonical controls.

\paragraph{2.6.5 Temporal holdout}\label{temporal-holdout}

To address label circularity \cite{Preibisch26}, IEDB epitopes
published after 2022 (post-2022, unseen by MHCflurry's
\textasciitilde pre-2021 training data) were compared against pre-2021
(in-distribution) epitopes at sensitivity threshold 0.95.

\paragraph{2.6.6 Evaluation}\label{evaluation}

Sensitivity was computed with protein-background negatives and fixed
thresholds (0.95 for the pan model; 500 nM for MHCflurry; 5th percentile
of canonical scores for MUNIS). Wilson 95\% confidence intervals are
reported for pooled estimates.

\subsection{3 Results}\label{results}

\subsubsection{3.1 Training data and architecture
comparison}\label{training-data-and-architecture-comparison}

MHCflurry labelling of 100,000 random 9-mers yielded 1.7\% SB, 5.1\% WB,
93.2\% NB; after balancing, 5,190 peptides (1,730/class). Across five
training seeds, the CNN achieved the highest accuracy (91.9\% $\pm$ 0.9\%;
Table 1, Figure 1), modestly ahead of the Deep FFN (90.6\% $\pm$ 1.2\%), the
FFN (89.8\% $\pm$ 1.1\%), and the LSTM (89.7\% $\pm$ 0.9\%) --- the top four
architectures within \textasciitilde2 percentage points --- while ResNet
(87.6\% $\pm$ 1.8\%) and Random Forest (79.7\% $\pm$ 0.4\%) were clearly worse.
MHCflurry itself scored 100\% on the test set it labelled --- an
expected ceiling, since the labels are its own predictions --- so the
meaningful comparison is the 8.1-point gap between the labeler's
decision boundary and the best finite-sample model that approximates it
(91.9\% for the CNN); independent-label circularity is instead
quantified by the 62.24\% NetMHCpan agreement (see Discussion). The WB
class was consistently the most difficult (F1 0.752--0.880).

\textbf{Table 1.} Model performance on the MHCflurry-labelled test set
(mean $\pm$ SD across five training seeds, n=519).

{\def\LTcaptype{none} % do not increment counter
\begin{longtable}[]{@{}lcc@{}}
\toprule\noalign{}
Model & Accuracy (\%) & Parameters \\
\midrule\noalign{}
\endhead
\bottomrule\noalign{}
\endlastfoot
CNN & \textbf{91.9 $\pm$ 0.9} & 1,053,863 \\
Deep FFN & 90.6 $\pm$ 1.2 & 152,823 \\
FFN (Jessen) & 89.8 $\pm$ 1.1 & 49,143 \\
LSTM & 89.7 $\pm$ 0.9 & 23,939 \\
ResNet & 87.6 $\pm$ 1.8 & 324,931 \\
Random Forest & 79.7 $\pm$ 0.4 & -- \\
\end{longtable}
}

\textdagger MHCflurry evaluated on its own labels scores 100\% --- an expected
ceiling, not evidence of superior performance; the 8.1-point gap to the
CNN reflects the information loss of approximating the labeler's
decision boundary from finite samples.

The modest spread among the top four architectures reflects the
positional nature of peptide--MHC binding: anchor positions p2/p9
dominate, so architectures that preserve positional information perform
comparably. ResNet and the Random Forest were clearly worse, consistent
with the absence of translation invariance in a 9$\times$20 encoding and, for
the Random Forest, the high dimensionality of the BLOSUM62 encoding.

\subsubsection{3.2 Label-source controls}\label{label-source-controls}

Label source, not architecture, sets the practical ceiling of what a
model can learn (Table 2, Figure 2). Training on MHCflurry labels
reached 90.6\% accuracy, whereas a deterministic PSSM label source
reached 82.6\% (on a separate 15,000-peptide pool, so not a clean
head-to-head comparison) and an i.i.d. random-label null fell to chance
(34.5 $\pm$ 1.7\%, close to the 33.3\% expected for three balanced classes).
The random-label result confirms the models learn genuine signal rather
than noise; the PSSM--MHCflurry gap, though confounded by pool size,
suggests MHCflurry's labels are at least as learnable as a hand-crafted
scoring matrix.

\textbf{Table 2.} Deep FFN performance across label sources (accuracy,
mean $\pm$ SD across five seeds where noted).

{\def\LTcaptype{none} % do not increment counter
\begin{longtable}[]{@{}lcl@{}}
\toprule\noalign{}
Label source & Accuracy (\%) & Notes \\
\midrule\noalign{}
\endhead
\bottomrule\noalign{}
\endlastfoot
DFFN on MHCflurry labels & 90.6 $\pm$ 1.2 & 5,190 peptides, 5 seeds \\
PSSM (biophysics) & 82.6 & separate 15,000-peptide pool, single seed \\
Random (i.i.d.) & 34.5 $\pm$ 1.7 & chance-level null, 5 seeds \\
\end{longtable}
}

\subsubsection{3.3 Encoding}\label{encoding}

On the FFN (the architecture used for all encoding comparisons),
BLOSUM62 generalized stably under cross-validation (92.2 $\pm$ 0.4\%;
per-fold 92.2/92.4/92.1/92.7/91.6\%). ESM-2 t6 per-position embeddings
scored 93.3\% on a single split but collapsed to 82.2 $\pm$ 1.0\% under CV
--- the 2,880-dimensional features overfit the 4,671-sample training set
(Table 3, Figure 7). Mean-pooled t6 embeddings collapsed to 65.9\%,
confirming that positional information is essential. The t12 variant
(4,320-dim) overfit similarly (90.9--91.5\%). The single-split 93.3\%
result reflects a favourable partition rather than a genuine encoding
advantage at this training scale; at the data scale typical of
single-allele studies, compact biologically motivated encodings such as
BLOSUM62 remain the more robust choice, and ESM-2's representational
advantage may require substantially larger training corpora to
materialise.

\textbf{Table 3.} Classification accuracy by encoding strategy.

{\def\LTcaptype{none} % do not increment counter
\begin{longtable}[]{@{}
  >{\raggedright\arraybackslash}p{(\linewidth - 6\tabcolsep) * \real{0.2857}}
  >{\centering\arraybackslash}p{(\linewidth - 6\tabcolsep) * \real{0.1429}}
  >{\centering\arraybackslash}p{(\linewidth - 6\tabcolsep) * \real{0.3714}}
  >{\raggedright\arraybackslash}p{(\linewidth - 6\tabcolsep) * \real{0.2000}}@{}}
\toprule\noalign{}
\begin{minipage}[b]{\linewidth}\raggedright
Encoding
\end{minipage} & \begin{minipage}[b]{\linewidth}\centering
Dim
\end{minipage} & \begin{minipage}[b]{\linewidth}\centering
Accuracy (\%)
\end{minipage} & \begin{minipage}[b]{\linewidth}\raggedright
Notes
\end{minipage} \\
\midrule\noalign{}
\endhead
\bottomrule\noalign{}
\endlastfoot
BLOSUM62 & 180 & 92.2 $\pm$ 0.4 (CV) & Stable CV generalisation \\
ESM-2 t6 mean pool & 320 & 65.9 & Positional info destroyed \\
ESM-2 t12 per-pos & 4,320 & 90.9 & Overfitting \\
ESM-2 t12 + L2 & 4,320 & 91.5 & Partially mitigated \\
ESM-2 t6 per-pos & 2,880 & 93.3 (82.2 $\pm$ 1.0 CV) & Single-split
optimistic \\
\end{longtable}
}

\subsubsection{3.4 Held-out validation and
translation}\label{held-out-validation-and-translation}

The Deep FFN achieved 95.8\% sensitivity (46/48), 75.0\% specificity
(15/20), ROC AUC 0.953 against IEDB epitopes (Table 5, Figure 3), after
excluding \texttt{ILRGSVAHK} (influenza NP 265--273), which is
HLA-B*08:01- rather than A*02:01-restricted. The two remaining false
negatives (\texttt{QYDPVAALF}, \texttt{TLGIVCPIC}) carried non-canonical
p9 anchors (F, C). An independent predictor, NetMHCpan-4.1 (queried
through the IEDB web API, which does not pin a reproducible tool
version), achieved 83.3\% sensitivity (40/48) on the same set. This gap
partly reflects label circularity (44/49 of the epitopes overlap
MHCflurry's training corpus), but NetMHCpan's pan-specific training on
MS-elution-based labels makes this an indicative rather than clean
comparison; the more direct circularity evidence is the 62.24\% label
agreement (323/519 test peptides) reported in the Discussion. To
contextualise the benchmark, we include a zero-parameter anchor-checker
baseline (a rule classifying any 9-mer with canonical p2/p9 residues as
a binder): it achieves sensitivity 1.0 on this canonical-motif-dominated
set, matching the trained model (a comparison with only 3 discordant
peptides, too low-powered for a formal test); the anchor rule has no
tunable score threshold, so its precision---and hence F1---is fixed by
the test-set composition rather than a free parameter --- confirming
that the canonical set overstates generalisation and that the
informative signal lies in the non-canonical/hard-tail sensitivity. The
five false positives were homopolymers ( $\leq$ 2 unique amino acids) with
canonical anchors --- biologically implausible binders whose low
sequence complexity the classifier fails to penalise. A domain-knowledge
filter ( $\leq$ 2 unique amino acids  $
ightarrow$  non-binder) removed all five, but we
report 75.0\% as the primary, unfiltered specificity to avoid post-hoc
inflation. To break this circularity, we additionally scored the Deep
FFN against a large label-independent set of experimentally measured
affinities (IEDB binding-affinity corpus, accessed 2026-08-02) --- 3,438 unique 9-mers with measured IC50 \textless500 nM
and 3,086 with IC50 \textgreater5,000 nM (the 500--5,000 nM grey zone
was excluded to yield unambiguous binders and non-binders), restricted
to direct binding assays (competitive fluorescence/radioactivity) and
excluding MS-elution-derived, predictor-dependent entries. The
MHCflurry-trained Deep FFN recognised 94.4\% of the experimental binders
and rejected 74.5\% of the non-binders (threshold 0.5; ROC AUC 0.942);
the rebalanced per-allele model reached 89.8\% sensitivity and 89.1\%
specificity (AUC 0.945). This benchmark validates the quality of the
training labels --- that MHCflurry's computational labels agree with
experimentally measured binding --- rather than generalisation to unseen
peptides, since the peptides are drawn from the same IEDB
binding-affinity corpus on which the labeler was trained; generalisation
to unseen peptides is instead established by the temporal holdout
(\S3.11). Protein scanning of 10 proteins (3,536 windows) recovered 10/11
known epitopes (91\%; Figure 4), including immunodominant
\texttt{GILGFVFTL} (M1), \texttt{NLVPMVATV} (pp65), and
\texttt{SLLMWITQC} (NY-ESO-1); the single miss was the MART-1 clinical
variant \texttt{ELAGIGILTV}, whose P2 leucine is a non-natural
anchor-modified residue. Mutation scanning of 16 hotspot mutations
identified 7 epitope-altering events (Table 4, Figure 5). p53 R248W
created the neoepitope \texttt{MNWRPILTI} (R $
ightarrow$ W at p3 introduces a bulky
hydrophobic residue improving groove fit); KRAS G12V created
\texttt{YKLVVVGAV} via G $
ightarrow$ V at p9 creating a canonical C-terminal anchor.
p53 R248W is among the most frequent TP53 missense mutations across
cancers (\textasciitilde3.5\% of somatic TP53 mutations, IARC), and
mutation-specific T-cell responses have been documented in patients and
in HLA-A2/DRB1-transgenic mice (Malekzadeh et al., 2020; Quandt et al.,
2018); KRAS G12C enhanced an existing epitope (WB $
ightarrow$ SB), suggesting
synergy with pharmacological G12C inhibition. Three mutations (KRAS
G13D, KRAS A146T, p53 R249S) destroyed existing epitopes, representing
potential immune-evasion mechanisms.

\textbf{Table 4.} Epitope-altering cancer hotspot mutations.

{\def\LTcaptype{none} % do not increment counter
\begin{longtable}[]{@{}llllc@{}}
\toprule\noalign{}
Mutation & WT peptide & Mutant peptide & Effect & $\Delta$ score \\
\midrule\noalign{}
\endhead
\bottomrule\noalign{}
\endlastfoot
p53 R248W & MNRRPILTI & MNWRPILTI & CREATED & +0.41 \\
KRAS G12V & YKLVVVGAG & YKLVVVGAV & CREATED & +0.48 \\
KRAS G12V & LVVVGAGGV & LVVVGAVGV & ENHANCED & +0.31 \\
KRAS G12C & LVVVGAGGV & LVVVGACGV & ENHANCED & +0.31 \\
KRAS G13D & GVGKSALTI & DVDKSALTI & DESTROYED & $-$0.47 \\
KRAS A146T & GIPFIETSA & GIPFIETST & DESTROYED & $-$0.42 \\
p53 R249S & GMNRRPILT & GMNRSPILT & DESTROYED & $-$0.16 \\
\end{longtable}
}

\textbf{Table 5.} IEDB benchmark performance.

{\def\LTcaptype{none} % do not increment counter
\begin{longtable}[]{@{}lc@{}}
\toprule\noalign{}
Metric & Value \\
\midrule\noalign{}
\endhead
\bottomrule\noalign{}
\endlastfoot
Sensitivity & 95.8\% (46/48) \\
NetMHCpan-4.1 sensitivity & 83.3\% (40/48) \\
Specificity & 75.0\% (15/20) \\
Precision & 90.2\% \\
F1 & 0.929 \\
ROC AUC & 0.953 \\
\end{longtable}
}

\subsubsection{3.5 A rebalanced pan model matches per-allele specialists
on four of five
alleles}\label{a-rebalanced-pan-model-matches-per-allele-specialists-on-four-of-five-alleles}

With protein-background negatives and a fixed threshold (n=25
non-canonical binders per allele), per-allele models reached
non-canonical sensitivity A*02:01 0.84, A*01:01 0.88, A*03:01 1.00,
B*07:02 0.96, B*44:03 0.96 (Table 6; Figure 8). The rebalanced pan model
matched this on four of five alleles (A*02:01 0.84, A*01:01 0.88,
A*03:01 1.00, B*44:03 1.00), but B*07:02 was a notable exception: the
pan model dropped to 0.60 versus 0.96 per-allele (0.64 in the 15-allele pan16 panel; Wilson 95\% CIs
{[}0.41, 0.77{]} vs {[}0.81, 0.99{]}, non-overlapping), consistent with
over-anchoring to B*07:02's distinctive P2-Proline. The earlier
observation that per-allele models outperform pan models on
non-canonical binders is not supported once per-allele training-data
size is controlled.

\textbf{Table 6. Non-canonical sensitivity by allele (protein-background
negatives, n=25 per allele, @0.95).}

{\def\LTcaptype{none} % do not increment counter
\begin{longtable}[]{@{}lcc@{}}
\toprule\noalign{}
Allele & Per-allele & Pan \\
\midrule\noalign{}
\endhead
\bottomrule\noalign{}
\endlastfoot
A*02:01 & 0.84 & 0.84 \\
A*01:01 & 0.88 & 0.88 \\
A*03:01 & 1.00 & 1.00 \\
B*07:02 & 0.96 & \textbf{0.60} \\
B*44:03 & 0.96 & 1.00 \\
\end{longtable}
}

\subsubsection{3.6 Canonical binders are recognized at
\textasciitilde0.98 by every
predictor}\label{canonical-binders-are-recognized-at-0.98-by-every-predictor}

On n = 11,225 canonical positives, all three predictors were saturated:
DFFN 0.979, MHCflurry 0.986, MUNIS 0.950 (Table 7). The canonical binder
problem is effectively solved by all current methods.

\subsubsection{3.7 Non-canonical (Type B) binders plateau at
\textasciitilde0.65 across predictor
classes}\label{non-canonical-type-b-binders-plateau-at-0.65-across-predictor-classes}

On n = 3,112 Type B positives, sensitivity was DFFN 0.607, MHCflurry
0.668, MUNIS 0.647 (Table 7; Figure 9). The three values fall within a
narrow band (0.607--0.668), far below the \textasciitilde0.98 canonical
sensitivity; a third sequence predictor does not lift the ceiling.

\textbf{Table 7. Predictor sensitivity by binder class (pooled, 5
alleles).}

{\def\LTcaptype{none} % do not increment counter
\begin{longtable}[]{@{}lccc@{}}
\toprule\noalign{}
Predictor & Canonical (n=11,225) & Type A (n=151) & Type B (n=3,112) \\
\midrule\noalign{}
\endhead
\bottomrule\noalign{}
\endlastfoot
DFFN (sequence) & 0.979 & 0.477 & 0.607 \\
MHCflurry (sequence) & 0.986 & 0.669 & 0.668 \\
MUNIS (sequence) & 0.950 & --- & 0.647 \\
\end{longtable}
}

\subsubsection{3.8 Type A (tolerated weak anchors) is inherently
hard}\label{type-a-tolerated-weak-anchors-is-inherently-hard}

On n = 151 Type A positives, DFFN achieved 0.477 and MHCflurry 0.669
(MUNIS was not evaluated on this small set). Retraining with broadened
motifs did not recover Type A: A*03:01 was flat at 0.634, and B*07:02
Type B sensitivity dropped from 0.680 to 0.341 (Table 8). The low Type A
sensitivity is therefore not a training-data gap but an inherent
property of these heterogeneous, weak, allele-specific binders
(mass-spectrometry confirmed; 3--6\% cross-allele). The canonical-motif
bias is what makes the model work.

\textbf{Table 8. Broad-motif retraining does not recover Type A.}

{\def\LTcaptype{none} % do not increment counter
\begin{longtable}[]{@{}lccc@{}}
\toprule\noalign{}
Allele & Type A before & Type A after & Type B after \\
\midrule\noalign{}
\endhead
\bottomrule\noalign{}
\endlastfoot
A*03:01 & --- & 0.634 & --- \\
A*02:01 & 0.08 & 0.24 & --- \\
B*07:02 & --- & --- & 0.341 (from 0.680) \\
\end{longtable}
}

\subsubsection{3.9 The extreme tail is prediction-resistant,
cross-predictor}\label{the-extreme-tail-is-prediction-resistant-cross-predictor}

The two false negatives from our single-allele IEDB benchmark
(QYDPVAALF, TLGIVCPIC) and the KRAS G12V P2-Lys epitope YKLVVVGAV were
scored as non-binders by all three predictors (MHCflurry 3.8--13.8  $\mu$ M;
DFFN 0.000--0.254; MUNIS 3/4 low), despite experimental evidence of
peptide-induced HLA-A*02:01-restricted cytotoxic T-cell responses ---
albeit the weakest of the three mutants screened in that study (Wu et
al., 2020). A single counterexample --- QYDPVAALF rescued by MUNIS
(0.972 vs.~sequence 0.000 / 6,199 nM) --- is encouraging but n = 1.

\subsubsection{3.10 Scaling: one model approximates a 15-allele
panel}\label{scaling-one-model-approximates-a-15-allele-panel}

The rebalanced pan16 reached pooled non-canonical sensitivity 0.819
(307/375, @0.95), with 12/15 alleles at 0.76--1.0 (Table 9; Figure 10).
The weak tail --- B*27:05 (0.08), B*51:01 (0.52), B*07:02 (0.64),
A*24:02 (0.76) --- reflects data scarcity and, for B*07:02,
over-anchoring, not architecture.

\textbf{Table 9. Pan16 non-canonical sensitivity by allele (15-allele
panel, @0.95).}

{\def\LTcaptype{none} % do not increment counter
\begin{longtable}[]{@{}lclc@{}}
\toprule\noalign{}
Allele & Sens & Allele & Sens \\
\midrule\noalign{}
\endhead
\bottomrule\noalign{}
\endlastfoot
A*02:01 & 0.88 & A*11:01 & 1.00 \\
A*01:01 & 0.80 & A*24:02 & 0.76 \\
A*03:01 & 1.00 & B*08:01 & 0.88 \\
B*07:02 & 0.64 & B*15:01 & 0.96 \\
B*44:03 & 1.00 & B*27:05 & 0.08 \\
& & B*35:01 & 0.84 \\
& & B*40:01 & 1.00 \\
& & B*49:01 & 1.00 \\
& & B*51:01 & 0.52 \\
& & B*57:01 & 0.92 \\
\textbf{Pooled} & \textbf{0.819} & & \\
\end{longtable}
}

\subsubsection{3.11 Temporal holdout breaks the label-circularity
concern}\label{temporal-holdout-breaks-the-label-circularity-concern}

Post-2022 epitopes (unseen by the labeler) were recognized at sens@0.95:
A*02:01 0.824, A*03:01 0.936, B*07:02 0.973 --- preserved versus
pre-2021 (0.879/0.905/0.898), with no systematic drop in the post-2022
era (Table 10; Figure 11). The model is not overfit to a single
training-data era.

\textbf{Table 10. Temporal holdout: pre-2021 vs.~post-2022 sensitivity
(@0.95).}

{\def\LTcaptype{none} % do not increment counter
\begin{longtable}[]{@{}lcc@{}}
\toprule\noalign{}
Allele & Pre-2021 (in-distribution) & Post-2022 (unseen) \\
\midrule\noalign{}
\endhead
\bottomrule\noalign{}
\endlastfoot
A*02:01 & 0.879 & 0.824 \\
A*03:01 & 0.905 & 0.936 \\
B*07:02 & 0.898 & 0.973 \\
\end{longtable}
}

\subsubsection{3.12 Learned feature
representations}\label{learned-feature-representations}

Feature extraction from the Deep FFN penultimate dense layer revealed
biologically meaningful structure (Figure 6). t-SNE visualisation of its
45-dimensional penultimate representation showed partial class
separation: SB peptides clustered distinctly (91\% of SB peptides fell
within the SB cluster), while WB peptides formed a continuous transition
zone between SB and NB clusters --- consistent with the graded affinity
boundary in the 50--500 nM range. Misclassified peptides (primarily
WB--NB boundary cases) mapped to the cluster overlap region, confirming
that representation quality degrades in the same ambiguity zone where
label contamination has the greatest impact. Integrated-gradients
analysis identified position-specific anchor contributions: p9 dominated
SB predictions, p2 and p3 contributed most to WB predictions, and NB
predictions were driven by the absence of anchor-characteristic
patterns. Hierarchical clustering separated validated epitopes into
coherent clusters, with false negatives forming a distinct branch
characterised by absent anchor-position activation. Five features showed
Pearson r \textgreater{} 0.7 with binding score; by a 1,000-iteration
permutation test, approximately 0.3 features would exceed this threshold
by chance, suggesting these represent genuine binding-relevant
representations rather than spurious correlations.

\subsection{4 Discussion}\label{discussion}

\subsubsection{Methodological warnings}\label{methodological-warnings}

Three apparent findings in this project dissolved under controlled
re-analysis, each illustrating how benchmark composition can manufacture
false conclusions. First, an apparent per-allele superiority over a pan
model was a 10.6$\times$ training-data-size artifact, not biology. Second, a
claimed `label quality dominates architecture' result collapsed once the
degenerate random-label null (which performed at chance,
\textasciitilde33\%) was removed --- the realistic label effect is
smaller than the architecture range. Third, a reported `simple
architecture beats complex ones' ranking was an artifact of stale model
checkpoints: re-training with fixed seeds showed a CNN instead modestly
outperforms the Deep FFN. We recommend that peptide-MHC benchmarks
report multi-seed estimates with confidence intervals,
protein-background negatives, frozen data snapshots, and a
state-of-the-art comparator such as NetMHCpan, to avoid these failure
modes. The same failure modes recur beyond peptide-MHC: data leakage
from protein language-model pretraining distorts downstream task
performance whenever train/test splits are not pretraining-aware
\cite{Hermann24}, and leakage is now recognised as a field-wide
reproducibility problem affecting hundreds of papers across seventeen
fields \cite{Kapoor23}. Dedicated guidance now exists for
avoiding leakage in biological machine learning \cite{Bernett24}
and for reporting clinical prediction models \cite{Collins24}; Box
1 instantiates these general principles for peptide-MHC benchmarking
specifically.

\textbf{Box 1. Benchmark hygiene recommendations.} The three artifacts
translate into reusable practices for immunoinformatics and
machine-learning benchmarks: (1) equalise training-data size across
model arms and report per-allele sample counts; (2) use a true i.i.d.
random-label null (chance  $\approx$  33\% for three balanced classes) rather than
a degenerate control; (3) freeze data snapshots, model checkpoints, and
label versions, and report multi-seed mean $\pm$ SD with confidence
intervals; (4) report label provenance and break label circularity with
a temporal holdout and an independent state-of-the-art comparator
(NetMHCpan). Each failure mode reproduced a false conclusion that was
caught only by controlled re-analysis.

Architecture choice has a modest but measurable effect on peptide-MHC
binding prediction. Across five seeds, a CNN (91.9\%) outperformed a
Deep FFN (90.6\%) by \textasciitilde1.3 percentage points (not a formal
significance test; the gap is within seed-to-seed variability), while a
Random Forest was clearly worse (79.7\%), for a maximum
inter-architecture range of \textasciitilde12 percentage points. The top
four architectures fell within \textasciitilde2 percentage points,
indicating that data quality and the decision boundary matter more than
architectural complexity for this task. Random-label controls performed
at chance (\textasciitilde33\%), confirming that the observed
performance reflects genuine signal rather than noise overfitting.

A PSSM-derived label source yielded lower accuracy than MHCflurry labels
(82.6\% vs 90.6\%, though on a separate peptide pool and therefore not a
clean comparison). This does not support the hypothesis that
deterministic biophysical labels are inherently superior to ML-derived
labels; rather, it suggests that MHCflurry's labels are at least as
learnable as a simple PSSM. We emphasise that training on any
computational label --- including MHCflurry --- re-deposits predictor
bias into downstream models, and recommend that benchmarks report label
provenance and the WB-class F1 as a standard diagnostic.

An independent state-of-the-art predictor, NetMHCpan-4.1, agreed with
the MHCflurry labels on only 62.24\% (323/519) of the same held-out test
peptides --- substantially below the 90.6\% achieved by a model trained
directly on those labels. This \textasciitilde28-point gap quantifies
the label-circularity limitation: our trained models learn MHCflurry's
specific decision boundary, whereas an independent predictor disagrees
on roughly one peptide in three, showing that the reported accuracies
measure fidelity to a particular labeler rather than true binding.

BLOSUM62 proved more robust than ESM-2 at single-allele training scales;
high-dimensional per-position embeddings overfit small datasets. The
2,880-dimensional per-position ESM-2 features create an unfavourable
feature-to-sample ratio for the 4,671-sample training set, and the
single-split 93.3\% result reflects a favourable partition rather than a
genuine encoding advantage. The shallower 6-layer ESM-2 variant
outperformed the deeper 12-layer model, possibly because t12's
additional transformer layers encode long-range protein-level
interactions that introduce noise at the short 9-mer peptide scale.
Architectural complexity beyond depth with batch normalisation gave
diminishing returns; a simple PSSM label source yielded lower accuracy
(82.6\%) than MHCflurry labels on its own pool, showing that deep
learning adds value over a hand-crafted scoring matrix for this task.
The language-model route may still be viable at single-allele scale if
the model is continued-pre-trained on HLA-associated peptides rather
than used as a frozen embedder, as such pre-training yields the largest
affinity gains precisely at small-data alleles \cite{Mares25}.

The multi-allele extension shows the recipe is portable: one rebalanced
pan model matches per-allele specialists on four of five alleles
(B*07:02 excepted) and reaches 0.819 pooled sensitivity on 15 alleles,
approximating NetMHCpan's practical value with a simple DFFN. This
aligns with the pan-specific paradigm: MixMHCpred v3.0's pan-allele
model now outperforms state-of-the-art MHC-I ligand predictors across
alleles and species \cite{Tadros25}. Yet non-canonical
(motif-failing) binders plateau at \textasciitilde0.65 sensitivity
across three independent sequence-based predictors: our Deep FFN
(0.607), MHCflurry (0.668), and the ESM-2-based MUNIS (0.647) fall
within a narrow band (0.61--0.67), far below the \textasciitilde0.98
canonical sensitivity, indicating this ceiling is shared across sequence
architectures of varying complexity rather than specific to our simple
network. The energy-compensated tail is prediction-resistant to all
three. The canonical binder problem is effectively solved; the
non-canonical tail defines the current frontier. This frontier extends
beyond anchor-motif violations to non-coding, fusion-derived, and
post-translationally modified peptides, for which current predictors
remain under-powered \cite{Vo25}.

\subsubsection{Translational
implications}\label{translational-implications}

The cancer neoepitope candidates identified here are testable
translational hypotheses. p53 R248W, which created \texttt{MNWRPILTI},
is among the most frequent TP53 missense mutations across human cancers
(\textasciitilde3.5\% of somatic TP53 mutations, IARC) with documented
gain-of-function oncogenic activity. A vaccine targeting this neoepitope
could, only if immunogenicity is experimentally confirmed, apply across
multiple tumour types; validation (e.g., IFN- $\gamma$  ELISpot using
HLA-A*02:01+ donor PBMCs) is required before any clinical development.
KRAS G12C, which enhanced an existing epitope (WB $
ightarrow$ SB), is the druggable
mutation targeted by sotorasib and adagrasib, suggesting a rational
combination immunotherapy strategy. The three epitope-destroying
mutations (KRAS G13D, KRAS A146T, p53 R249S) represent potential
immune-evasion mechanisms. We note that binding prediction alone does
not establish immunogenicity: prior work shows that only a small
fraction of top-predicted MHC-I binders are confirmed immunogenic, and
that peptide--MHC stability predicts immunogenicity better than affinity
(Harndahl et al., 2012; Backert and Kohlbacher, 2015). High-throughput
thermostability profiling now quantifies pMHC stability at
immunopeptidome scale and improves neoepitope immunogenicity prediction
\cite{Jappe20}, and transfer learning further improves
kinetic-stability prediction from limited data \cite{Fasoulis24}.
The binding-to-immunogenicity gap remains the central bottleneck in
translating computational epitope prediction to clinical vaccine design.

\subsubsection{Limitations}\label{limitations}

Training labels are MHCflurry-derived rather than experimental
affinities; we quantified the circularity by cross-referencing the 49
curated IEDB epitopes against MHCflurry's training corpus --- 44/49
(89.8\%) overlapped, a figure computed on all 49 before the exclusion of
ILRGSVAHK (HLA-B*08:01) that reduces the evaluated set to 48; all five
non-overlapping epitopes were nonetheless correctly classified (5/5,
100\%), providing limited evidence of generalisation beyond the training
distribution. The single-allele design was deliberate --- holding the
allele constant isolates label-source variation from allele-driven
confounds --- but the multi-allele extension addresses the
generalisability question directly. The multi-allele estimates rest on
IEDB-derived positives; experimental (non-IEDB) validation of the
extreme tail is outstanding, and a pMHC-fine-tuned structure model
(beyond generic AlphaFold-Multimer) might lift the Type B plateau. All
translational predictions and the cancer neoepitope candidates require
experimental confirmation (e.g., IFN- $\gamma$  ELISpot or tetramer staining).
Antigen-processing steps (proteasomal cleavage, TAP transport) are not
modelled.

\subsection{5 Conclusion}\label{conclusion}

Architecture choice has a modest effect on peptide-MHC binding
prediction: a CNN modestly outperformed a Deep FFN, the top four
architectures fell within \textasciitilde2 percentage points, and a
Random Forest was clearly worse. Random-label controls confirmed
chance-level performance, ruling out noise overfitting. Canonical
binding is largely solved by the anchor motif --- a zero-parameter
anchor rule matches the trained model --- while non-canonical binders
plateau at \textasciitilde0.65 sensitivity across three independent
sequence-based predictors, defining the current frontier. Beyond
peptide-MHC, the three artifacts we corrected --- data-size imbalance, a
degenerate null, and stale checkpoints --- illustrate a general
benchmark-composition failure mode that reproduces false conclusions
across machine-learning benchmarks; the checklist in Box 1 is offered as
a reusable safeguard. Experimental validation of the non-canonical tail
is the decisive next step.

\section*{Figures}
\begin{figure}[H]
\centering
\includegraphics[width=0.85\textwidth]{figures/Figure1.pdf}
\caption{Model performance comparison across six architectures on the MHCflurry-labelled test set (accuracy, macro F1, WB F1).}
\end{figure}

\begin{figure}[H]
\centering
\includegraphics[width=0.85\textwidth]{figures/Figure2.pdf}
\caption{Architecture comparison (Table 1) and label-source controls: MHCflurry (90.6\%), PSSM (82.6\%), random labels at chance ($\sim$33\%).}
\end{figure}

\begin{figure}[H]
\centering
\includegraphics[width=0.85\textwidth]{figures/Figure3.pdf}
\caption{IEDB benchmark: ROC curve (AUC 0.953), confusion matrix, per-epitope scores.}
\end{figure}

\begin{figure}[H]
\centering
\includegraphics[width=0.85\textwidth]{figures/Figure4.pdf}
\caption{Protein epitope scanning: SB/WB/NB distribution across 10 proteins and top epitope candidates.}
\end{figure}

\begin{figure}[H]
\centering
\includegraphics[width=0.85\textwidth]{figures/Figure5.pdf}
\caption{Cancer hotspot mutation analysis: delta scores and wild-type/mutant pairs for 7 epitope-altering mutations.}
\end{figure}

\begin{figure}[H]
\centering
\includegraphics[width=0.85\textwidth]{figures/Figure6.pdf}
\caption{t-SNE of the Deep FFN 45-dimensional penultimate-layer representation, coloured by binding class (SB/WB/NB).}
\end{figure}

\begin{figure}[H]
\centering
\includegraphics[width=0.85\textwidth]{figures/Figure7.pdf}
\caption{ESM-2 embedding comparison across five encoding strategies.}
\end{figure}

\begin{figure}[H]
\centering
\includegraphics[width=0.85\textwidth]{figures/Figure8.pdf}
\caption{Rebalanced pan model non-canonical sensitivity (5-seed mean $\pm$ SD) across five common HLA-I alleles. Per-allele values are tabulated in Table 6 rather than overlaid, as the per-allele benchmark (n=25 broad-motif) is not directly comparable to the large non-canonical set.}
\end{figure}

\begin{figure}[H]
\centering
\includegraphics[width=0.85\textwidth]{figures/Figure9.pdf}
\caption{Predictor plateau: canonical versus Type B (non-canonical) sensitivity for Deep FFN, MHCflurry, and MUNIS (sequence); dashed line marks the $\sim$0.65 Type B plateau.}
\end{figure}

\begin{figure}[H]
\centering
\includegraphics[width=0.85\textwidth]{figures/Figure10.pdf}
\caption{Pan16: per-allele non-canonical sensitivity for the 15-allele panel, sorted ascending; dashed line marks the pooled 0.819.}
\end{figure}

\begin{figure}[H]
\centering
\includegraphics[width=0.85\textwidth]{figures/Figure11.pdf}
\caption{Temporal holdout: non-canonical sensitivity at threshold 0.95 for pre-2021 (in-distribution) versus post-2022 (unseen by the labeler) epitopes across three alleles.}
\end{figure}

\subsection{Data and code
availability}\label{data-and-code-availability}

All code, trained models (per-allele \texttt{deep\_ffn\_HLA-A\_*.keras}
and pan
\texttt{deep\_ffn\_pan\_seed*.keras}/\texttt{deep\_ffn\_pan16.keras}),
feature matrices, prediction results, the curated multi-allele dataset
(\texttt{type\_ab\_dataset.json}, 14,488 positives), and the
experimental-label benchmark (6,524 measured IC50 values) are available
at github.com/woodhaha/peptide-mhc-benchmark. The per-allele Deep FFN is
distributed as the \texttt{pmhcbenchmark} Python package, with a
\texttt{pmhc-predict} command-line tool and a standalone Windows GUI
(\texttt{pmhc\_predict.exe}); exported weights
(\texttt{pmhc\_models.npz}) enable tensorflow-free inference. The
single-allele benchmark-hygiene analysis and the multi-allele
pan-vs-per-allele benchmark (\S3.5--3.11) together constitute the
complete study.

\subsection{Author contributions}\label{author-contributions}

Conceptualization: Yunshi Huang, Weisheng Chen, Zhuha Zhou. Methodology:
Yunshi Huang, Weisheng Chen. Formal analysis: Yunshi Huang, Weisheng
Chen, Yiqi Cai. Investigation: Yunshi Huang, Weisheng Chen, Yiqi Cai.
Validation: Xiaopeng Zhuang, Zhuxian Zhou. Writing -- original draft:
Yunshi Huang, Weisheng Chen, Yiqi Cai. Writing -- review \& editing:
Xiaopeng Zhuang, Zhuxian Zhou, Xiangyang Xue, Zhuha Zhou. Supervision:
Xiangyang Xue, Zhuha Zhou.

\subsection{Competing interests}\label{competing-interests}

The authors declare no competing interests.

\subsection{Funding}\label{funding}

This research did not receive specific grant funding.

\subsection{Ethics statement}\label{ethics-statement}

This study used only publicly available, de-identified data from the
Immune Epitope Database (IEDB) and public protein databases; no human
subjects, animals, or patient-derived samples were involved, and no
ethics approval or informed consent was required.

\bibliographystyle{plos2015}
\bibliography{manuscript_plos}

\end{document}